\documentclass[11pt, aspectratio=169]{beamer}
\usepackage[style=authoryear-icomp, backend=biber]{biblatex}
\usetheme{jcar}
\usepackage{booktabs}

%\includeonlyframes{a} % Para compilar solo en la que se está trabajando

%\addbibresource{refs.bib}

\begin{document}
\author[Juan C. Arboleda Rivera]{Juan Camilo Arboleda Rivera}
\title{Pruebas de hipótesis}
%\subtitle{}
%\titlegraphic{\includegraphics[width=2cm]{udea}}
\institute[U. de A.]{Instituto de Biología\\
  Facultad de Ciencias Exactas y Naturales\\
  Universidad de Antioquia\\
  juan.arboleda2@udea.edu.co}
\date{v: \today}
%\subject{}

\begin{frame}[plain]
  \maketitle
\end{frame}

\begin{frame}
 \frametitle{Contenidos}
 \tableofcontents
\end{frame}

\section{Introducción}

\begin{frame}{¿Por qué necesitamos pruebas de hipótesis?}

Supongamos que estudiamos una especie de planta.

\begin{block}{Pregunta biológica}
¿La aplicación de un fertilizante aumenta el crecimiento de las plantas?
\end{block}

Observamos:
\vspace{-1em}
\[
\bar{x}_{fertilizante}=18.4\text{ cm}
\]
\vspace{-1em}
\[
\bar{x}_{control}=16.9\text{ cm}
\]

Hay una diferencia de:

\[
18.4-16.9=1.5\text{ cm}
\]

\pause

Pero...

\begin{alertblock}{Pregunta estadística}
¿Una diferencia de 1.5 cm es suficientemente grande como para ser difícil
de explicar por variación aleatoria?
\end{alertblock}

\end{frame}

% ------------------------------------------------------------
\begin{frame}{Variabilidad: una característica de los datos biológicos}

La biología presenta variabilidad debido a:

\begin{itemize}[<+->]
    \item Diferencias genéticas
    \item Edad y estado fisiológico
    \item Ambiente
    \item Error de medición
    \item Heterogeneidad espacial
    \item Variación temporal
    \item Factores no controlados
\end{itemize}

\vspace{0.3cm}

Por eso:

\vspace{-.5em}

\[
\text{diferencia observada}
=
\text{señal}
+
\text{variación}
\]

\pause

La estadística intenta cuantificar cuánto de la diferencia observada
es compatible con la variabilidad esperada.

\end{frame}

% ------------------------------------------------------------
\begin{frame}{La lógica de una prueba de hipótesis}

\begin{enumerate}[<+->]
    \item Planteamos una hipótesis sobre la población.
    \item Recolectamos una muestra.
    \item Calculamos un estadístico (un número del cual se conoce su
    \emph{distribución} y que me permite estimar la probabilidad de obtener
    un determinado resultado por azar).
    \item Determinamos qué tan compatible es el estadístico con una
    hipótesis nula ($H_0$).
    \item Tomamos una decisión basada en una regla previamente definida.
\end{enumerate}

\vspace{0.4cm}

\begin{center}
\Large
\[
\boxed{
H_0
\longrightarrow
\text{datos}
\longrightarrow
\text{estadístico}
\longrightarrow
p
}
\]
\end{center}

\end{frame}

% ------------------------------------------------------------
\begin{frame}{Hipótesis nula y alternativa}

\begin{block}{Hipótesis nula ($H_0$)}
Representa generalmente una ausencia de diferencia, asociación o efecto
bajo el modelo planteado.
\end{block}

\begin{block}{Hipótesis alternativa ($H_1$ o $H_A$)}
Representa el patrón que queremos evaluar como evidencia frente a $H_0$.
\end{block}

\vspace{0.3cm}

Ejemplo:

\[
H_0:\mu_{fert}=\mu_{control}
\]

\[
H_A:\mu_{fert}\neq\mu_{control}
\]

\end{frame}

% ------------------------------------------------------------
%\begin{frame}{Hipótesis unilaterales y bilaterales}
%
%%### Bilateral
%\vspace{-1em}
%\[
%H_A:\mu\neq\mu_0
%\]
%
%Pregunta:
%
%\begin{quote}
%¿La media es diferente?
%\end{quote}
%
%%### Unilateral
%
%\[
%H_A:\mu>\mu_0
%\]
%
%Pregunta:
%
%\begin{quote}
%¿La media es mayor?
%\end{quote}
%
%También:
%\vspace{-1em}
%\[
%H_A:\mu<\mu_0
%\]
%
%\begin{alertblock}{Importante}
%La dirección de la hipótesis debe estar justificada por la pregunta
%científica
%y definida antes de observar los resultados.
%\end{alertblock}
%
%\end{frame}

% ------------------------------------------------------------
\begin{frame}{Ejemplo: efecto de un contaminante}

Queremos estudiar si una exposición a un contaminante afecta
la tasa fotosintética.

Variable:

\[
Y = \text{tasa fotosintética}
\]

Pregunta:

\begin{quote}
¿Las plantas expuestas presentan una tasa fotosintética diferente
a las plantas control?
\end{quote}

Entonces:

\[
H_0:\mu_E=\mu_C
\]

\[
H_A:\mu_E\neq\mu_C
\]

\end{frame}

% ------------------------------------------------------------
\begin{frame}
%[label=a]
{La hipótesis biológica vs. Hipótesis estadística}

\centering

La hipótesis biológica está relacionada con la hipótesis estadística pero
\textbf{NO} son iguales.\\
Cada prueba estadística tiene su Hipótesis nula
definida.

\pause

\begin{center}
\begin{tabular}{p{0.65\textwidth} p{0.3\textwidth}}
\toprule
\textbf{Pregunta biológica} & \textbf{Hipótesis estadística}\\
\midrule
¿Determinado tratamiento cambia el peso de un insecto? &
$\mu_T \neq \mu_C$\\[0.2cm]

¿El fertilizante aumenta el crecimiento? &
$\mu_F > \mu_C$\\[0.2cm]

¿Dos hábitats tienen distinta abundancia? &
$\mu_1 \neq \mu_2$\\[0.2cm]

¿La temperatura está asociada con el número de huevos? &
$\rho \neq 0$\\
\bottomrule
\end{tabular}
\end{center}

\end{frame}

\section{Estadístico de prueba}

% ------------------------------------------------------------
\begin{frame}{El estadístico de prueba}

Una prueba transforma los datos en un número que mide la discrepancia
entre lo observado y lo esperado bajo $H_0$.

Ejemplo: prueba $t$

\[
t=
\frac{\bar{x}-\mu_0}
{s/\sqrt{n}}
\]

donde:

\begin{itemize}
    \item $\bar{x}$ = media observada,
    \item $\mu_0$ = media bajo $H_0$,
    \item $s$ = desviación estándar,
    \item $n$ = tamaño de muestra.
\end{itemize}

\end{frame}

% ------------------------------------------------------------
\begin{frame}{¿Qué significa el estadístico?}

Supongamos:

\[
\bar{x}=12.4
\]
\vspace{-.5em}
\[
\mu_0=10
\]
\vspace{-.5em}
\[
s=4
\]

\[
n=16
\]

Entonces:

\[
t=
\frac{12.4-10}{4/\sqrt{16}}
=
\frac{2.4}{1}
=
2.4
\]

La diferencia observada equivale a aproximadamente
$2.4$ errores estándar respecto al valor planteado por $H_0$.

\end{frame}

\begin{frame}
[label=a]
{La distribución del estadístico de prueba es conocida}{}
 \centering
 \textbf{Ejemplo:}\\
 La distribución t de Student\\
 \includegraphics[width=.8\linewidth]{t-de-Student.png}
\end{frame}

\begin{frame}
[label=a]
{}{}
 \centering
 \includegraphics[width=\linewidth]{t-de-Student.jpg}
\end{frame}

% ------------------------------------------------------------
\begin{frame}{¿Qué es el valor-$p$?}

El valor-$p$ responde a una pregunta específica:

\begin{block}{}
Si $H_0$ fuera cierta y el modelo estadístico fuera apropiado,
¿qué tan compatible sería observar un resultado como el obtenido
(o más extremo)?
\end{block}

\vspace{0.4cm}

Un valor-$p$ pequeño indica que los datos son poco compatibles con
$H_0$ bajo el modelo utilizado.

\end{frame}

% ------------------------------------------------------------
\begin{frame}{Lo que un valor-$p$ NO significa}

Si:

\[
p=0.03
\]

\textbf{NO significa:}

\begin{itemize}
    \item ``Hay 3\% de probabilidad de que $H_0$ sea cierta''.
    \item ``Hay 97\% de probabilidad de que $H_A$ sea cierta''.
    \item ``El efecto es grande''.
    \item ``El resultado es biológicamente importante''.
    \item ``El experimento está probado''.
\end{itemize}

\pause

\begin{alertblock}{Recordatorio}
El valor-$p$ se calcula suponiendo un modelo bajo $H_0$.
No es una probabilidad de que una hipótesis sea verdadera.
\end{alertblock}

\end{frame}

% ------------------------------------------------------------
\begin{frame}{Nivel de significancia $\alpha$}

Antes de analizar los resultados podemos definir:

\[
\alpha=0.05
\]

Regla convencional:

\[
p<\alpha
\quad\Rightarrow\quad
\text{rechazar }H_0
\]

\[
p\geq\alpha
\quad\Rightarrow\quad
\text{no rechazar }H_0
\]

\begin{alertblock}{Lenguaje importante}
``No rechazar $H_0$'' no significa ``demostrar que $H_0$ es verdadera''.
\end{alertblock}

\end{frame}

\section{Errores tipo I y tipo II}

% ------------------------------------------------------------
\begin{frame}{Errores tipo I y tipo II}

Tenemos dos posibles errores:

\begin{center}
\begin{tabular}{c|cc}
& $H_0$ verdadera & $H_0$ falsa\\
\hline
Rechazar $H_0$ &
\textbf{Error tipo I} &
Decisión correcta\\
No rechazar $H_0$ &
Decisión correcta &
\textbf{Error tipo II}
\end{tabular}
\end{center}

\vspace{0.3cm}

\[
P(\text{Error tipo I})=\alpha
\]

\[
P(\text{Error tipo II})=\beta
\]

\end{frame}

% ------------------------------------------------------------
\begin{frame}{Ejemplo biológico: falso positivo}

Supongamos que evaluamos un nuevo tratamiento contra una enfermedad vegetal.

\[
H_0:\text{el tratamiento no tiene efecto}
\]

Un error tipo I ocurre si:

\[
\text{rechazamos }H_0
\]

cuando realmente el tratamiento no tiene efecto.

\begin{block}{Interpretación}
Concluir que existe un efecto cuando, bajo el modelo, realmente
no existe.
\end{block}

\end{frame}
% ------------------------------------------------------------
\begin{frame}{Ejemplo biológico: falso negativo}

Ahora supongamos que el tratamiento realmente funciona.

Pero nuestro estudio no encuentra evidencia suficiente para rechazar $H_0$.

Esto corresponde a un:

\[
\boxed{\text{Error tipo II}}
\]

Puede ocurrir cuando:

\begin{itemize}
    \item $n$ es pequeño,
    \item existe mucha variabilidad,
    \item el efecto es pequeño,
    \item las mediciones son poco precisas.
\end{itemize}

\end{frame}

\begin{frame}
%[label=a]
{}{}
 \centering
 \includegraphics[height=\textheight]{type-i-and-type-ii-errors.jpg}
\end{frame}

% ------------------------------------------------------------
\begin{frame}{Potencia estadística}

La potencia es:

\[
\boxed{\text{Potencia}=1-\beta}
\]

Representa la probabilidad de detectar un efecto de una magnitud
determinada cuando ese efecto realmente existe, bajo las condiciones
del diseño y del análisis.

La potencia aumenta generalmente con:

\begin{itemize}
    \item mayor tamaño de muestra,
    \item mayor tamaño del efecto,
    \item menor variabilidad,
    \item mediciones más precisas.
\end{itemize}

\end{frame}


% ------------------------------------------------------------
\begin{frame}{Ejemplo conceptual de potencia}

Supongamos que queremos detectar una diferencia de:

\[
\Delta=2\text{ cm}
\]

entre dos tratamientos.

\begin{itemize}
    \item Con $n=5$ por grupo puede ser difícil.
    \item Con $n=20$ por grupo tenemos más información.
    \item Con $n=100$ por grupo podemos detectar diferencias pequeñas.
\end{itemize}

\begin{alertblock}{Idea clave}
Un resultado no significativo puede deberse a ausencia de efecto,
pero también a un estudio con poca capacidad para detectarlo.
\end{alertblock}

\end{frame}

% ------------------------------------------------------------
\begin{frame}{Tamaño del efecto}

La significancia estadística no mide directamente la magnitud del efecto.

Ejemplo:

\[
\Delta = 0.2\text{ cm}
\]

puede ser estadísticamente significativo con una muestra enorme.

Mientras que:

\[
\Delta = 5\text{ cm}
\]

puede no ser significativo con una muestra muy pequeña.

\begin{block}{En biología}
Debemos reportar tanto la evidencia estadística como la magnitud
y la incertidumbre del efecto.
\end{block}

\end{frame}

\section{Ejemplos de pruebas estadísticas}

% ------------------------------------------------------------
\begin{frame}{Ejemplo 1: una muestra}

Pregunta:

\begin{quote}
¿La longitud promedio de las hojas de una población es diferente
de 10 cm?
\end{quote}

Hipótesis:

\[
H_0:\mu=10
\]

\[
H_A:\mu\neq10
\]

Si las condiciones son razonables:

\[
\text{prueba }t\text{ de una muestra}
\]

\end{frame}

% ------------------------------------------------------------
\begin{frame}{Ejemplo 2: dos grupos independientes}

Tenemos dos grupos:

\begin{itemize}
    \item Plantas control.
    \item Plantas tratadas.
\end{itemize}

Variable respuesta:

\[
Y=\text{biomasa seca}
\]

Pregunta:

\[
\mu_T-\mu_C\neq0?
\]

Una opción frecuente:

\[
\boxed{\text{prueba }t\text{ para dos muestras}}
\]

\end{frame}

% ------------------------------------------------------------
\begin{frame}{Prueba $t$: idea intuitiva}

La prueba compara:

\[
\text{diferencia entre medias}
\]

con:

\[
\text{variabilidad esperada de esa diferencia}
\]

En términos intuitivos:

\[
t=
\frac{\text{señal}}
{\text{ruido}}
\]

\pause

Cuanto mayor sea la diferencia respecto a su incertidumbre,
más evidencia puede existir contra $H_0$.

\end{frame}

% ------------------------------------------------------------
\begin{frame}{Ejemplo 3: datos pareados}

Medimos la presión arterial de animales:

\[
\text{Antes}
\qquad\text{y}\qquad
\text{Después}
\]

Los datos están relacionados porque:
\vspace{-.5em}
\[
\boxed{\text{cada animal aporta dos mediciones}}
\]

Podemos trabajar con las diferencias:

\[
D_i=Después_i-Antes_i
\]

y probar:

\[
H_0:\mu_D=0
\]

Se usa:
\[
\boxed{\text{prueba }t\text{ pareada}}
\]


\end{frame}

% ------------------------------------------------------------
\begin{frame}{Ejemplo 4: más de dos grupos}

Queremos comparar la biomasa de plantas bajo:

\begin{itemize}
    \item Control
    \item Fertilizante A
    \item Fertilizante B
    \item Fertilizante C
\end{itemize}

Una prueba ANOVA puede evaluar:

\[
H_0:
\mu_C=\mu_A=\mu_B=\mu_C
\]

frente a:

\[
H_A:\text{al menos una media difiere}
\]

\end{frame}

%\appendix
%\AtNextBibliography{\footnotesize}
%\begin{frame}[plain, allowframebreaks]{Bibliografía}{}
% \bibitemsep 1em
% \printbibliography[heading=none]
%\end{frame}
\end{document}

%%%%%%%
Ejemplos:

Pregunta biológica: ¿un fertilizante aumenta el crecimiento de las plantas?

Se puede proporcionar a cada estudiante una muestra simulada:

Planteamiento:

    H0:μfertilizante−μcontrol=0

    H1:μfertilizante−μcontrol≠0

Simulación

En lugar de comenzar con una prueba t, los estudiantes pueden:

1. Calcular la diferencia observada entre los grupos.
2. Mezclar todas las observaciones.
3. Redistribuir aleatoriamente las plantas entre dos grupos.
4. Calcular nuevamente la diferencia.
5. Repetir 1.000 veces.
6. Construir la distribución de diferencias esperadas bajo H0.
7. Comparar la diferencia observada con esa distribución.

Así se visualiza que la prueba de hipótesis pregunta qué tan compatible es
el resultado observado con la hipótesis nula, no simplemente si "los
promedios son diferentes".

%%%%%%%%%%%%%%%

Depredación y coloración de escarabajos

Esta puede ser especialmente buena para biología evolutiva.

Imagine una población de escarabajos:

    50 % verdes

    50 % marrones

Los estudiantes simulan que las aves capturan escarabajos al azar.

Después de 100 capturas observan, por ejemplo:

    38 verdes

    62 marrones

Pregunta:

    ¿Esta diferencia podría aparecer simplemente por azar si las aves
    realmente no prefieren ningún color?

Se simulan muchas muestras de 100 escarabajos bajo H0:p=0.5.

Los estudiantes pueden visualizar:

H₀ → simulación → distribución de resultados → resultado observado →
evidencia contra H₀

Es una forma muy natural de conectar estadística con selección natural.

%%%%%%%%%%%%%%%%%%
Sí. Para estudiantes de biología suele funcionar mejor que la prueba de
hipótesis aparezca como una **simulación de un experimento biológico**, y no
como una fórmula aislada.

 Aquí tienes varias ideas que pueden convertirse en actividades de clase:

 ## 1\. Monedas → semillas germinadas

 **Objetivo:** introducir la idea de distribución bajo $H_0$, variabilidad
 aleatoria y valor p.

 - Supongamos que, en condiciones normales, **50 % de las semillas
 %germinan**.
- $H_0: p=0.50$
- $H_1: p>0.50$
- Los estudiantes simulan 20, 50 o 100 semillas mediante lanzamientos de
moneda.
- Repiten el experimento muchas veces y observan cuántas simulaciones
producen, por azar, resultados como 70 %, 75 %, etc.
- Después se introduce la pregunta central:

 > **Si $H_0$ fuera cierta, ¿qué tan sorprendente sería obtener nuestro
 resultado?**

 Esto permite introducir el **valor p** de manera intuitiva antes de
 calcularlo.

---

 ## 2\. Efecto de un fertilizante sobre el crecimiento

 **Pregunta biológica:** ¿un fertilizante aumenta el crecimiento de las
 plantas?

 Se puede proporcionar a cada estudiante una muestra simulada:

 | Control (cm) | Fertilizante (cm) |
| --- | --- |
| 12.1 | 14.2 |
| 11.8 | 13.9 |
| 13.0 | 14.5 |
| 12.5 | 13.7 |
| ... | ... |

Planteamiento:

 - $H_0:\mu_{\text{fertilizante}}-\mu_{\text{control}}=0$
- $H_1:\mu_{\text{fertilizante}}-\mu_{\text{control}}\neq0$

 ### Simulación

 En lugar de comenzar con una prueba _t_, los estudiantes pueden:

 1. Calcular la diferencia observada entre los grupos.
2. Mezclar todas las observaciones.
3. Redistribuir aleatoriamente las plantas entre dos grupos.
4. Calcular nuevamente la diferencia.
5. Repetir 1.000 veces.
6. Construir la distribución de diferencias esperadas bajo $H_0$.
7. Comparar la diferencia observada con esa distribución.

 Así se visualiza que **la prueba de hipótesis pregunta qué tan compatible
 es el resultado observado con la hipótesis nula**, no simplemente si "los
 promedios son diferentes".

---

 ## 3\. Depredación y coloración de escarabajos

 Esta puede ser especialmente buena para biología evolutiva.

 Imagine una población de escarabajos:

 - 50 % verdes
- 50 % marrones

 Los estudiantes simulan que las aves capturan escarabajos al azar.

 Después de 100 capturas observan, por ejemplo:

 - 38 verdes
- 62 marrones

 Pregunta:

 > ¿Esta diferencia podría aparecer simplemente por azar si las aves
 realmente no prefieren ningún color?

 Se simulan muchas muestras de 100 escarabajos bajo $H_0:p=0.5$.

 Los estudiantes pueden visualizar:

 **H₀ → simulación → distribución de resultados → resultado observado →
 evidencia contra H₀**

 Es una forma muy natural de conectar estadística con **selección natural**.

---

 ## 4\. Frecuencia de un alelo

 Para genética de poblaciones:

 Supongamos que la frecuencia esperada del alelo $A$ es:

 $$
p=0.30
$$

 Se simulan muestras de 100 individuos y se cuenta cuántos presentan el
 alelo.

 Luego se pregunta:

 > Si la frecuencia verdadera fuera 0.30, ¿qué frecuencias muestrales
 esperaríamos observar?

 Los estudiantes generan, por ejemplo, **10.000 poblaciones simuladas** y
 construyen un histograma.

 Después se coloca una línea vertical en el valor observado.

 Esto permite introducir:

 - distribución muestral;
- variabilidad;
- hipótesis nula;
- región extrema;
- valor p;
- error tipo I.

---

 ## 5\. "Detectives": ¿el tratamiento realmente funciona?

 Una actividad interesante es proporcionar a cada grupo una **base de datos
 diferente**, pero decirles únicamente:

 > "Un investigador afirma que el tratamiento aumenta la supervivencia."

 Cada grupo recibe datos de, por ejemplo, 30 individuos tratados y 30
 controles.

 Los estudiantes deben decidir:

 1. ¿Cuál es la variable respuesta?
2. ¿Cuál sería $H_0$?
3. ¿Cuál sería $H_1$?
4. ¿Qué patrón esperaríamos si $H_0$ fuera cierta?
5. ¿Qué tan extremo es nuestro resultado?
6. ¿Qué conclusión estadística podemos sacar?

 Después se comparan las decisiones entre grupos.

 Esto ayuda a evitar que la enseñanza quede reducida a **"calcular un p y
 compararlo con 0,05"**.

---

 ## Una simulación particularmente potente para una primera clase

 Yo estructuraría una actividad alrededor de esta pregunta:

 > **¿Puede aparecer una diferencia aparentemente grande simplemente por
 azar?**

 Por ejemplo, se plantea:

 **Experimento:** 20 plantas reciben tratamiento A y 20 tratamiento B.

 Se obtiene:

 - A: media = 12,4 cm
- B: media = 10,8 cm
- diferencia = 1,6 cm

 Antes de realizar cualquier prueba estadística, los estudiantes tienen que
 apostar intuitivamente:

 - "La diferencia es real".
- "Podría ser azar".
- "No hay suficiente información".

 Luego comienza la simulación:

 **Supongamos que el tratamiento no tiene ningún efecto.**

 Se mezclan las 40 mediciones y se redistribuyen aleatoriamente en dos
 grupos de 20.

 Cada redistribución produce una diferencia.

 Después de 1.000–10.000 redistribuciones aparece una distribución:

```
                 distribución bajo H₀

                    ▄▄▄
                  ▄█████▄
                ▄█████████▄
              ▄████████████▄
          ▄▄▄████████████████▄▄▄
       ▄█████████████████████████▄
───────┴───────────────────────────┴───────
     -2       -1       0       1       2
                              ↑
                         diferencia
                          observada
```

 La pregunta deja de ser:

 > "¿Cuál es la fórmula de la prueba?"

 y pasa a ser:

 > **"¿Qué tan frecuente sería obtener una diferencia de 1,6 cm o mayor si
 realmente no hubiera efecto?"**

 Ese cambio conceptual suele ser mucho más importante que memorizar la
 fórmula.

 ### Secuencia didáctica que puedes repetir

 **Pregunta biológica → hipótesis nula → simulación bajo H₀ → distribución
 nula → resultado observado → valor p → conclusión**

 Y después puedes conectar la simulación con las pruebas formales:

 - una media → prueba _t_;
- dos grupos → _t_ de Student/Welch;
- proporciones → prueba binomial o χ²;
- frecuencias genotípicas → χ²;
- asociación entre variables → correlación/regresión;
- varios grupos → ANOVA.

 Si quieres, también puedo prepararte una **simulación interactiva lista
 para usar en clase**, donde los estudiantes puedan cambiar el tamaño de
 muestra, el efecto biológico y el nivel de significancia y ver cómo cambia
 la distribución nula, el valor _p_ y los errores tipo I/II.